% =============================================================================
% A Bayesian Statistical Framework for Detecting Engine-Assisted Play
% in Online Gomoku Competitions
% =============================================================================
\documentclass[11pt,a4paper]{article}

% --- Packages ---
\usepackage[utf8]{inputenc}
\usepackage[T1]{fontenc}
\usepackage{lmodern}
\usepackage{amsmath,amssymb,amsthm}
\usepackage{graphicx}
\usepackage{booktabs}
\usepackage{hyperref}
\usepackage[margin=2.5cm]{geometry}
\usepackage{algorithm}
\usepackage{algpseudocode}
\usepackage{enumitem}
\usepackage{caption}
\usepackage{subcaption}
\usepackage{xcolor}
\usepackage{float}
\usepackage{natbib}

% --- Theorem environments ---
\newtheorem{definition}{Definition}
\newtheorem{proposition}{Proposition}
\newtheorem{theorem}{Theorem}
\newtheorem{lemma}{Lemma}
\newtheorem{remark}{Remark}

% --- Custom commands ---
\newcommand{\E}{\mathbb{E}}
\newcommand{\Prob}{\mathbb{P}}
\newcommand{\R}{\mathbb{R}}
\DeclareMathOperator*{\argmax}{arg\,max}
\DeclareMathOperator*{\argmin}{arg\,min}

% =============================================================================
\title{
    A Bayesian Statistical Framework for Detecting Engine-Assisted Play\\
    in Online Gomoku Competitions
}

\author{
    % TODO: Replace with actual author information
    Author Name \\
    \textit{Affiliation} \\
    \texttt{email@example.com}
}

\date{\today}

% =============================================================================
\begin{document}

\maketitle

% =============================================================================
% ABSTRACT
% =============================================================================
\begin{abstract}
The integrity of online board game competitions is increasingly threatened by
players who covertly use artificial intelligence engines during live matches.
This paper presents a rigorous Bayesian statistical framework for detecting
engine-assisted play in Gomoku (Five-in-a-Row), a combinatorial strategy game.
Our approach models the \emph{winrate delta}---the difference between the
engine-optimal move evaluation and the player's actual move evaluation---as a
random variable drawn from parametric distributions that differ between
legitimate human players and engine-assisted players. We develop the framework
within a general family of non-negative distributions (Exponential, Gamma,
Weibull), and introduce a \emph{two-component mixture model} to capture the
bimodal nature of human play, with a Dirichlet Process extension for
automatic component selection. A \emph{context-aware complexity metric} weights
evidence by positional difficulty, and a \emph{tempered likelihood} mechanism
prevents trivial positions from distorting posterior estimates. The framework
performs sequential Bayesian updates over the course of a game, integrating
a Boltzmann choice model, information-theoretic divergence measures, and
temporal switch-point detection to identify mid-game engine activation. We
further introduce a formal likelihood ratio test grounded in Neyman--Pearson
theory for false-positive rate control. We describe the complete mathematical
formulation, the system architecture for engine integration, and the
classification decision boundaries. The proposed method provides a principled,
transparent, and mathematically grounded alternative to ad-hoc statistical
thresholds commonly used in current anti-cheating systems.

\medskip
\noindent
\textbf{Keywords:} Anti-cheating, Bayesian inference, Gomoku, engine-assisted
play detection, online gaming integrity, mixture model, Gamma distribution,
Weibull distribution, Dirichlet process, information theory, temporal analysis.
\end{abstract}

% =============================================================================
% 1. INTRODUCTION
% =============================================================================
\section{Introduction}
\label{sec:introduction}

Online board game platforms have experienced significant growth, enabling
players worldwide to compete in real-time tournaments. However, the remote and
unsupervised nature of online play creates an avenue for a specific form of
cheating: covertly consulting an AI engine during a live match. In Gomoku
(also known as Five-in-a-Row or Renju in its professional variant), modern
engines achieve superhuman playing strength, making engine-assisted play both
highly effective and difficult to distinguish from expert human play through
casual observation alone.

Existing anti-cheating measures in online gaming have largely relied on
behavioral heuristics---such as timing patterns, mouse movement tracking, or
tab-switching detection---which are easily circumvented by sophisticated
cheaters. Statistical approaches have been explored in chess
\citep{regan2011intrinsic}, but relatively little formal work has been published
for Gomoku and its variants. Furthermore, many deployed systems rely on ad-hoc
thresholds (e.g., ``if accuracy exceeds 90\%, flag the player'') without a
coherent probabilistic foundation, leading to both false positives against
strong legitimate players and false negatives against selective cheaters.

This paper addresses these limitations by developing a complete Bayesian
statistical framework that:

\begin{enumerate}[label=(\roman*)]
    \item Models human and cheater move-quality distributions using parametric
          families grounded in statistical theory.
    \item Incorporates \emph{positional context} through a complexity metric
          that modulates the evidential weight of each move.
    \item Performs online (sequential) Bayesian updates with tempered
          likelihoods, ensuring that trivial positions do not corrupt the
          posterior estimate.
    \item Employs information-theoretic divergence measures and temporal
          change-point detection to identify partial or mid-game engine use.
\end{enumerate}

\subsection{Contributions}

The main contributions of this work are:

\begin{enumerate}
    \item \textbf{Winrate Delta Model.} We formalize the concept of
          \emph{winrate delta} and establish that the Exponential distribution
          family provides a natural parametric model for move quality under both
          hypotheses.
    \item \textbf{Two-Component Mixture Model.} We propose a mixture of
          Exponentials for the human hypothesis that accounts for the bimodal
          nature of human play (careful moves interspersed with occasional
          blunders).
    \item \textbf{Context-Aware Complexity.} We derive a composite complexity
          metric from candidate move evaluations that weights Bayesian evidence
          by positional difficulty, preventing trivial positions from generating
          misleading signals.
    \item \textbf{Tempered Bayesian Updates.} We introduce a tempering mechanism
          where the likelihood exponent equals the positional complexity,
          providing a mathematically principled method for evidence weighting.
    \item \textbf{Temporal Switch-Point Detection.} We develop a temporal model
          that detects abrupt changes in play quality, identifying possible
          mid-game engine activation.
    \item \textbf{Complete System Design.} We describe the full architecture for
          integrating an external game engine as an evaluation oracle within the
          detection pipeline.
\end{enumerate}

\subsection{Paper Organization}

Section~\ref{sec:related_work} reviews related work. Section~\ref{sec:framework}
presents the mathematical framework in detail. Section~\ref{sec:complexity}
develops the context-aware complexity model. Section~\ref{sec:bayesian}
describes the Bayesian inference engine. Section~\ref{sec:information_theory}
introduces the information-theoretic metrics. Section~\ref{sec:temporal}
covers temporal analysis and switch-point detection.
Section~\ref{sec:architecture} outlines the system architecture.
Section~\ref{sec:classification} defines the classification decision
boundaries. Section~\ref{sec:discussion} discusses strengths, limitations, and
ethical considerations. Section~\ref{sec:conclusion} concludes the paper.


% =============================================================================
% 2. RELATED WORK
% =============================================================================
\section{Related Work}
\label{sec:related_work}

\subsection{Engine-Assisted Play Detection in Chess}

The most extensively studied domain for engine-assisted play detection is chess.
Regan and Haworth~\citep{regan2011intrinsic} proposed an intrinsic performance
rating system that estimates a player's Elo rating from their moves alone,
comparing observed play quality to engine evaluations. Their approach uses a
logistic model to predict the probability that a player of a given rating would
select each candidate move. Significant deviations between the predicted rating
and the player's registered rating are flagged as suspicious.

The \emph{Let's Check} system used by ChessBase and similar platforms
\citep{guid2006computer} compares player moves against deep engine analysis,
computing aggregate statistics such as average centipawn loss and correlation
with engine top choices. While effective as a screening tool, these systems
typically rely on fixed thresholds rather than probabilistic inference.

\subsection{Statistical Methods for Fair Play}

Broader statistical approaches to fair play detection include hypothesis testing
frameworks~\citep{berger1987testing}, where the null hypothesis is honest play
and the alternative is engine assistance. The challenge lies in controlling the
false positive rate---strong human players naturally exhibit high move quality,
and a framework that flags top-level play as suspicious undermines competitive
integrity.

\subsection{Gomoku-Specific Considerations}

Gomoku differs from chess in several important respects that affect anti-cheating
analysis. The game has a narrower decision tree in many positions due to forced
tactical sequences (threats, double-threats), meaning that many ``perfect''
moves are also the obvious human choice. A robust detection framework must
therefore distinguish between positions where engine-optimal play is trivially
findable and positions where it requires deep calculation. This motivates our
development of the context-aware complexity model (Section~\ref{sec:complexity}).


% =============================================================================
% 3. MATHEMATICAL FRAMEWORK
% =============================================================================
\section{Mathematical Framework: The Winrate Delta Model}
\label{sec:framework}

\subsection{Definitions}

We assume access to an external game engine that, given a board position~$s$,
produces a set of $k$ candidate moves $\{m_1, m_2, \ldots, m_k\}$ along with
their associated \emph{winrate evaluations} $\{w_1, w_2, \ldots, w_k\}$, where
$w_i \in [0, 1]$ represents the probability of winning from the resulting
position if move~$m_i$ is played.

\begin{definition}[Winrate Delta]
\label{def:delta}
For a given position, let $w^* = \max_i w_i$ denote the winrate of the
engine's best candidate move, and let $w_p$ denote the winrate of the move
actually played by the player. The \emph{winrate delta} is defined as:
\begin{equation}
    \Delta = w^* - w_p, \quad \Delta \geq 0.
    \label{eq:delta}
\end{equation}
The non-negativity constraint is enforced by construction: the played move
cannot have a higher evaluation than the engine's best candidate (within the
set of evaluated moves). In the rare case where the played move was not among
the candidates, $\Delta$ is set to $\max(0, w^* - w_p)$.
\end{definition}

The winrate delta $\Delta$ captures the \emph{quality loss} incurred by the
player's move choice. A $\Delta$ of zero indicates the player found the
engine-optimal move; a large $\Delta$ indicates a suboptimal choice.

\begin{remark}
Throughout this paper, all winrates are expressed from the perspective of the
player whose moves are being analyzed. If the engine reports winrates from the
opponent's perspective (as is common in Gomoku protocols), the inversion
$w_{\text{player}} = 1 - w_{\text{opponent}}$ must be applied before computing
$\Delta$.
\end{remark}

\subsection{Distributional Family for the Winrate Delta}
\label{sec:dist_family}

Because $\Delta \geq 0$ by definition, the appropriate parametric models
belong to the family of non-negative continuous distributions. We consider
three candidates in order of increasing flexibility, and develop our baseline
framework using the Exponential while providing the Gamma and Weibull
generalizations for empirical model selection.

\subsubsection{Exponential Distribution (Baseline)}

The simplest model assumes $\Delta \sim \text{Exponential}(\lambda)$:
\begin{equation}
    f(\Delta \mid \lambda) = \lambda \, e^{-\lambda \Delta}, \quad \Delta \geq 0.
    \label{eq:exponential_pdf}
\end{equation}

The Exponential distribution offers parsimony (a single parameter $\lambda$)
and a direct encoding of the central tendency: $\E[\Delta] = 1/\lambda$.
Its memoryless property---$P(\Delta > t_1 + t_2 \mid \Delta > t_1) =
P(\Delta > t_2)$---implies that the magnitude of a suboptimal choice does
not depend on previous errors conditional on the current position.

However, this memoryless assumption may be violated in practice.
Players experiencing psychological momentum (``tilt'' or ``flow state'') may
exhibit serial correlation in their error magnitudes. We therefore present
two generalizations that relax this assumption.

\subsubsection{Gamma Distribution (Shape-Rate Generalization)}

The Gamma distribution generalizes the Exponential by introducing a shape
parameter $\alpha > 0$:
\begin{equation}
    f(\Delta \mid \alpha, \beta) = \frac{\beta^{\alpha}\,
    \Delta^{\alpha - 1}\, e^{-\beta\Delta}}{\Gamma(\alpha)},
    \quad \Delta \geq 0.
    \label{eq:gamma_pdf}
\end{equation}

When $\alpha = 1$, the Gamma reduces to $\text{Exponential}(\beta)$.
For $\alpha > 1$, the density is unimodal with mode at
$(\alpha - 1)/\beta$, enabling it to model delta distributions that
concentrate around a typical error magnitude rather than being monotonically
decreasing---a more realistic characterization of average-strength human
players whose errors cluster around a characteristic value.

The mean and variance are $\E[\Delta] = \alpha/\beta$ and
$\text{Var}(\Delta) = \alpha/\beta^2$, allowing independent control over
location and spread.

\subsubsection{Weibull Distribution (Tail Flexibility)}

The Weibull distribution offers direct control over tail behavior:
\begin{equation}
    f(\Delta \mid k, \lambda) = \frac{k}{\lambda}
    \left(\frac{\Delta}{\lambda}\right)^{k-1}
    \exp\!\left[-\left(\frac{\Delta}{\lambda}\right)^{k}\right],
    \quad \Delta \geq 0.
    \label{eq:weibull_pdf}
\end{equation}

The shape parameter $k$ controls the tail weight:
\begin{itemize}
    \item $k = 1$: reduces to $\text{Exponential}(1/\lambda)$.
    \item $k < 1$: heavier tail than Exponential (more frequent large
          errors)---suitable for human players prone to occasional severe
          blunders.
    \item $k > 1$: lighter tail (errors concentrate tightly)---suitable for
          cheater modeling where large deviations are extremely rare.
\end{itemize}

The Weibull is particularly relevant because human blunders are characterized
by rare but disproportionately large $\Delta$ values (heavy-tailed behavior),
while cheater errors are tightly concentrated (light-tailed behavior).

\subsubsection{Model Selection via Information Criteria}
\label{sec:model_selection}

Given observed data, the appropriate distributional family is selected using
the \emph{Akaike Information Criterion} (AIC) and \emph{Bayesian Information
Criterion} (BIC):
\begin{align}
    \text{AIC} &= 2k - 2\ln \hat{L}, \label{eq:aic}\\
    \text{BIC} &= k \ln n - 2\ln \hat{L}, \label{eq:bic}
\end{align}
where $k$ is the number of parameters, $n$ is the sample size, and
$\hat{L}$ is the maximized likelihood. Lower values indicate better
model--data fit after penalizing complexity. AIC favors predictive accuracy
while BIC favors parsimonious models. Additionally,
\emph{goodness-of-fit tests} (Kolmogorov--Smirnov or Anderson--Darling) should
be applied to verify that the chosen distribution adequately fits the
empirical delta distribution before deploying the system.

In the remainder of this paper, we develop the framework using the
Exponential family for notational clarity, with the understanding that
the Gamma or Weibull can be substituted when empirical evidence warrants it.

\subsection{Hypothesis Structure}

Under the Exponential baseline, the two hypotheses are characterized by
different rate parameters:
\begin{itemize}
    \item \textbf{Human hypothesis} ($H_0$): $\lambda_{\text{human}}$ is
          relatively small (e.g., $\lambda_{\text{human}} = 5$, corresponding to
          $\E[\Delta] = 0.20$), reflecting larger average winrate losses.
    \item \textbf{Cheater hypothesis} ($H_1$): $\lambda_{\text{cheater}}$ is
          large (e.g., $\lambda_{\text{cheater}} = 50$, corresponding to
          $\E[\Delta] = 0.02$), reflecting near-optimal play with very small
          average winrate losses.
\end{itemize}

Under the Gamma generalization, each hypothesis is parameterized by
$(\alpha_j, \beta_j)$ for $j \in \{H_0, H_1\}$, estimated via
maximum likelihood on labeled reference data.

\subsection{Maximum Likelihood Estimation}

Given a sequence of observed deltas $\Delta_1, \Delta_2, \ldots, \Delta_n$,
the MLE for the Exponential rate parameter is:
\begin{equation}
    \hat{\lambda}_{\text{MLE}} = \frac{n}{\sum_{i=1}^{n} \Delta_i}
    = \frac{1}{\bar{\Delta}}.
    \label{eq:lambda_mle}
\end{equation}

To handle the case $\Delta_i = 0$ (which occurs when the player finds the
exact engine-optimal move), we apply a floor $\epsilon = 0.001$:
\begin{equation}
    \Delta_i^{\text{adj}} = \max(\Delta_i, \, \epsilon).
    \label{eq:delta_adj}
\end{equation}

For the Gamma distribution, the MLE for $(\alpha, \beta)$ requires
numerical optimization of the log-likelihood:
\begin{equation}
    \ell(\alpha,\beta) = n\alpha\log\beta - n\log\Gamma(\alpha)
    + (\alpha-1)\sum_{i}\log\Delta_i - \beta\sum_{i}\Delta_i.
    \label{eq:gamma_mle}
\end{equation}

We define the \emph{Accuracy Index} as:
\begin{equation}
    \text{AccuracyIndex} = \frac{1}{\hat{\lambda}_{\text{MLE}}} = \bar{\Delta}^{\text{adj}},
    \label{eq:accuracy_index}
\end{equation}
which directly measures the average winrate loss. Lower values indicate play
closer to engine-optimal.


% =============================================================================
% 3.1 MIXTURE MODEL
% =============================================================================
\subsection{Two-Component Exponential Mixture for Human Play}
\label{sec:mixture}

A single Exponential distribution inadequately captures human play. In
practice, human players exhibit a bimodal pattern: the majority of their moves
are reasonably sound (small $\Delta$), but occasional moves are outright
blunders (large $\Delta$). A single Exponential with a low $\lambda$ would
overestimate the frequency of moderate errors while underweighting both
near-optimal play and severe blunders.

We therefore model the human delta distribution as a two-component Exponential
mixture:
\begin{equation}
    P(\Delta \mid H_0) = \pi \cdot \lambda_g \, e^{-\lambda_g \Delta}
    + (1 - \pi) \cdot \lambda_b \, e^{-\lambda_b \Delta},
    \label{eq:mixture_pdf}
\end{equation}
where:
\begin{itemize}
    \item $\pi \in (0, 1)$ is the mixing weight for the ``good play'' component
          (typically $\pi = 0.75$).
    \item $\lambda_g$ is the rate for good plays (e.g., $\lambda_g = 20$,
          corresponding to $\E[\Delta_g] = 0.05$).
    \item $\lambda_b$ is the rate for blunders (e.g., $\lambda_b = 3$,
          corresponding to $\E[\Delta_b] = 0.33$).
\end{itemize}

The log-likelihood for a single observation under the mixture model is computed
using the log-sum-exp trick for numerical stability:
\begin{align}
    a &= \log \pi + \log \lambda_g - \lambda_g \Delta, \label{eq:log_comp_good}\\
    b &= \log(1 - \pi) + \log \lambda_b - \lambda_b \Delta, \label{eq:log_comp_blunder}\\
    \log P(\Delta \mid H_0) &= \log\!\left(e^a + e^b\right)
    = m + \log\!\left(e^{a-m} + e^{b-m}\right),
    \label{eq:log_mixture}
\end{align}
where $m = \max(a, b)$ ensures that the dominant term does not cause numerical
overflow.

The cheater hypothesis retains a unimodal model because engine-assisted play
is consistently near-optimal with no blunder component. Under the Exponential
baseline:
\begin{equation}
    P(\Delta \mid H_1) = \lambda_c \, e^{-\lambda_c \Delta}.
    \label{eq:cheater_pdf}
\end{equation}

Under the Gamma generalization, the cheater hypothesis becomes
$\Delta \mid H_1 \sim \text{Gamma}(\alpha_c, \beta_c)$ with $\alpha_c$
typically near~1 (preserving near-Exponential shape) and $\beta_c \gg 1$
(tight concentration near zero).

\subsection{Parameter Estimation via Expectation-Maximization}
\label{sec:em}

The mixture parameters $\Theta = (\pi, \lambda_g, \lambda_b)$ should
ideally be estimated from labeled reference data via the
\emph{Expectation-Maximization} (EM) algorithm~\citep{dempster1977maximum}
rather than assigned heuristically. The EM algorithm iterates between:

\textbf{E-step:} Compute the responsibility of each component for each
observation~$\Delta_i$:
\begin{equation}
    \gamma_{i,1} = \frac{\pi \cdot \lambda_g e^{-\lambda_g \Delta_i}}
    {\pi \cdot \lambda_g e^{-\lambda_g \Delta_i}
    + (1-\pi) \cdot \lambda_b e^{-\lambda_b \Delta_i}}, \quad
    \gamma_{i,2} = 1 - \gamma_{i,1}.
    \label{eq:e_step}
\end{equation}

\textbf{M-step:} Update the parameters:
\begin{align}
    \pi^{\text{new}} &= \frac{1}{n}\sum_{i=1}^{n} \gamma_{i,1},
    \label{eq:m_step_pi}\\
    \lambda_g^{\text{new}} &= \frac{\sum_i \gamma_{i,1}}
    {\sum_i \gamma_{i,1} \Delta_i}, \quad
    \lambda_b^{\text{new}} = \frac{\sum_i \gamma_{i,2}}
    {\sum_i \gamma_{i,2} \Delta_i}.
    \label{eq:m_step_lambda}
\end{align}

The EM algorithm is guaranteed to increase the log-likelihood at each
iteration and converges to a local maximum~\citep{bishop2006pattern}.

\subsection{Dirichlet Process Mixture Model}
\label{sec:dpm}

The fixed two-component mixture (Section~\ref{sec:mixture}) assumes a priori
that human behavior decomposes into exactly two modes. In reality, human play
may exhibit additional modes (e.g., opening preparation, tactical
calculation, endgame technique), and the optimal number of components is
unknown.

The \emph{Dirichlet Process Mixture Model}
(DPMM)~\citep{ferguson1973bayesian, neal2000markov} provides a
nonparametric Bayesian solution that allows the number of components to be
inferred from the data:
\begin{equation}
    G \sim \text{DP}(\alpha_0, G_0), \qquad
    \theta_i \mid G \sim G, \qquad
    \Delta_i \mid \theta_i \sim F(\cdot \mid \theta_i),
    \label{eq:dpmm}
\end{equation}
where $\text{DP}(\alpha_0, G_0)$ is a Dirichlet Process with concentration
parameter $\alpha_0$ and base measure $G_0$ (e.g., a Gamma prior over
Exponential rate parameters), and $F(\cdot \mid \theta_i)$ is the component
density.

The DPMM marginalizes over all possible numbers of components
$K \in \{1, 2, 3, \ldots\}$, with the effective number of occupied clusters
governed by $\alpha_0$. This avoids the model selection problem entirely:
the data determines whether two, three, or more modes are needed.
Inference proceeds via Gibbs sampling or variational
methods~\citep{blei2006variational}.

While the DPMM is more computationally demanding than the fixed mixture,
it represents the principled extension for production systems with access to
large labeled datasets.


% =============================================================================
% 3.2 SOFTMAX CHOICE MODEL
% =============================================================================
\subsection{Boltzmann Choice Model}
\label{sec:softmax}

Complementing the delta-based analysis, we model the player's move
\emph{selection process} using a Boltzmann (softmax) distribution over
candidate winrates~\citep{luce1959individual}:
\begin{equation}
    P(m_i \mid \{w_j\}, \tau) = \frac{\exp(w_i / \tau)}{\sum_{j=1}^{k}
    \exp(w_j / \tau)},
    \label{eq:softmax}
\end{equation}
where $\tau > 0$ is a \emph{temperature} parameter. A low temperature
($\tau \to 0$) concentrates probability mass on the highest-winrate move
(engine-like behavior), while a high temperature ($\tau \to \infty$) produces
a uniform distribution over candidates (random play).

The temperature $\tau$ is estimated via grid-search maximum likelihood over
a range $[\tau_{\min}, \tau_{\max}]$:
\begin{equation}
    \hat{\tau} = \argmax_{\tau} \sum_{i=1}^{n} \log P(m_{c_i} \mid
    \{w_j^{(i)}\}, \tau),
    \label{eq:tau_mle}
\end{equation}
where $c_i$ is the index of the move chosen at position~$i$.

\subsubsection{Integration into the Classification Pipeline}

Unlike the delta model, the Boltzmann temperature captures the player's
\emph{selection strategy} across all candidates, not just the gap between
best and played. The estimated temperature is integrated into the final
classification as an auxiliary score:
\begin{equation}
    S_{\tau} = \frac{1}{\hat{\tau}},
    \label{eq:tau_score}
\end{equation}
where a high $S_{\tau}$ (low temperature) indicates engine-like selection
behavior. This score enters the classification through a logistic
aggregation layer (Section~\ref{sec:classification}).


% =============================================================================
% 4. CONTEXT-AWARE COMPLEXITY
% =============================================================================
\section{Context-Aware Position Complexity}
\label{sec:complexity}

Not all positions are equally informative. A position where only one move
is reasonable provides no evidence about the player's strength or engine use.
A position with multiple viable alternatives of substantially different quality
is highly informative.

We define a composite \emph{complexity score} $C \in [0, 1]$ from four
factors derived from the engine's candidate evaluations.

\subsection{Delta Vector}

Given $k$ candidate winrates $\{w_1, \ldots, w_k\}$ and the player's previous
winrate $w_{\text{prev}}$ (the evaluation after the opponent's last move),
define the \emph{delta vector}:
\begin{equation}
    d_i = w_i - w_{\text{prev}}, \quad i = 1, \ldots, k.
    \label{eq:delta_vector}
\end{equation}

The key metrics derived from this vector are:
\begin{align}
    d^* &= \max_i d_i, &  \bar{d} &= \frac{1}{k}\sum_{i=1}^{k} d_i,  &
    S &= d^* - \min_i d_i.
    \label{eq:delta_metrics}
\end{align}

\subsection{Complexity Factors}

\begin{definition}[Impact Factor]
\label{def:impact}
Measures how much the best available move changes the evaluation:
\begin{equation}
    F_{\text{impact}} = \min\!\left(\frac{|d^*|}{0.25},\; 1.0\right).
    \label{eq:impact}
\end{equation}
\end{definition}

\begin{definition}[Variance Factor]
\label{def:variance}
Measures how different the candidate evaluations are from each other:
\begin{equation}
    F_{\text{variance}} = \min\!\left(\frac{S}{0.40},\; 1.0\right).
    \label{eq:variance}
\end{equation}
\end{definition}

\begin{definition}[Criticality Factor]
\label{def:criticality}
Measures how uniquely important the best move is:
\begin{equation}
    F_{\text{criticality}} = \begin{cases}
        \displaystyle\frac{d^* - \bar{d}}{S} & \text{if } S > 0.01, \\[6pt]
        0 & \text{otherwise}.
    \end{cases}
    \label{eq:criticality}
\end{equation}
Clamped to $[0, 1]$.
\end{definition}

\begin{definition}[Phase Factor]
\label{def:phase}
Encodes game-phase relevance based on the best achievable winrate
$w^* = \max_i w_i$:
\begin{equation}
    F_{\text{phase}} = \begin{cases}
        0.3 & \text{if } w^* > 0.85 \quad\text{(winning)}, \\
        0.2 & \text{if } w^* < 0.15 \quad\text{(losing)}, \\
        0.5 + 0.5 \cdot (1 - 2|w^* - 0.5|) & \text{otherwise (balanced)}.
    \end{cases}
    \label{eq:phase}
\end{equation}
\end{definition}

\subsection{Combined Complexity Score}

The overall complexity is a weighted combination:
\begin{equation}
    C = 0.15 \cdot F_{\text{impact}} + 0.40 \cdot F_{\text{variance}}
    + 0.20 \cdot F_{\text{criticality}} + 0.25 \cdot F_{\text{phase}},
    \label{eq:complexity}
\end{equation}
with the variance factor receiving the largest weight, reflecting the principle
that decision difficulty is principally determined by the spread of candidate
evaluations.

\begin{proposition}[Trivial Position Override]
\label{prop:trivial}
If the position is trivially won ($w^* > 0.95$) or trivially lost
($w^* < 0.05$) and the variance factor is below~$0.05$, then:
\begin{equation}
    C = 0.10.
    \label{eq:trivial}
\end{equation}
This override ensures that positions with no meaningful decision contribute
minimal evidence to the Bayesian posterior.
\end{proposition}

\subsection{Context-Aware Accuracy}

We define a position-adjusted accuracy score for each move:
\begin{equation}
    A(\Delta, C) = \begin{cases}
        1.0 & \text{if } \Delta \leq 0.02, \\[4pt]
        \exp\!\left(-5 \cdot \Delta \cdot (1 - 0.3 \cdot C)\right) &
        \text{otherwise}.
    \end{cases}
    \label{eq:accuracy}
\end{equation}

The complexity modulation $(1 - 0.3C)$ ensures that errors in complex positions
are penalized less than errors in simple positions, reflecting the greater
difficulty of finding the best move.

\subsubsection{Strategic Bonus}

If the player's move creates high complexity for the opponent (measured by
opponent candidate variance), a strategic bonus of up to 15\% is applied:
\begin{equation}
    A_{\text{final}} = A \cdot (1 + 0.15 \cdot C_{\text{opp}}),
    \quad \text{if } A > 0.4 \text{ and } C_{\text{opp}} > 0.1,
    \label{eq:strategic_bonus}
\end{equation}
clamped to $[0, 1]$. This prevents the penalization of strategically rich
moves that sacrifice marginal winrate for positional complexity.


% =============================================================================
% 5. BAYESIAN INFERENCE
% =============================================================================
\section{Bayesian Inference Engine}
\label{sec:bayesian}

\subsection{Posterior Computation}

Given a sequence of winrate deltas $D = \{\Delta_1, \ldots, \Delta_n\}$, the
posterior probability of cheating is computed via Bayes' theorem:
\begin{equation}
    P(H_1 \mid D) = \frac{P(D \mid H_1) \cdot P(H_1)}{P(D \mid H_1) \cdot
    P(H_1) + P(D \mid H_0) \cdot P(H_0)},
    \label{eq:bayes}
\end{equation}
where $P(H_1)$ is the prior probability of cheating (set conservatively to
$0.01$) and $P(H_0) = 1 - P(H_1)$.

All computations are performed in log space to avoid numerical underflow.
Defining:
\begin{align}
    L_0 &= \log P(D \mid H_0) = \sum_{i=1}^{n} \log P(\Delta_i \mid H_0),
    \label{eq:log_lik_human}\\
    L_1 &= \log P(D \mid H_1) = \sum_{i=1}^{n} \log P(\Delta_i \mid H_1),
    \label{eq:log_lik_cheater}
\end{align}
the normalized posterior is:
\begin{align}
    \ell_0 &= L_0 + \log P(H_0), &
    \ell_1 &= L_1 + \log P(H_1), \nonumber\\
    m &= \max(\ell_0, \ell_1), \label{eq:logsumexp}\\
    P(H_1 \mid D) &= \frac{e^{\ell_1 - m}}{e^{\ell_0 - m} + e^{\ell_1 - m}}.
    \label{eq:posterior_normalized}
\end{align}

\subsection{Online (Sequential) Bayesian Updates}

For real-time analysis during a game in progress, we perform \emph{sequential}
Bayesian updates. After observing each new move with delta $\Delta_{n+1}$ and
complexity $C_{n+1}$, the posterior is updated as:
\begin{equation}
    P(H_1 \mid D_{1:n+1}) \propto P(\Delta_{n+1} \mid H_1)^{C_{n+1}} \cdot
    P(H_1 \mid D_{1:n}),
    \label{eq:sequential_update}
\end{equation}
where the current posterior serves as the prior for the next update.

\subsection{Tempered Likelihood}
\label{sec:tempered}

The key innovation is the use of the positional complexity $C$ as a
\emph{tempering exponent} on the likelihood. Instead of multiplying the
log-likelihood by a complexity weight (which lacks a probabilistic
interpretation), we raise the likelihood to the power of $C$:
\begin{equation}
    P(\Delta \mid H)^{\alpha} \quad \Longleftrightarrow \quad
    \alpha \cdot \log P(\Delta \mid H), \quad \alpha = C.
    \label{eq:tempering}
\end{equation}

This approach has the following mathematical properties:
\begin{enumerate}
    \item When $\alpha = 1$ (maximum complexity), the observation contributes
          full evidence.
    \item When $\alpha = 0$ (trivial position), the tempered likelihood equals
          $1$ (i.e., $\log P^0 = 0$), contributing zero evidence and leaving
          the posterior unchanged.
    \item For intermediate values, the evidence contribution scales smoothly
          and continuously.
    \item Tempered likelihoods are well-studied in the statistical literature
          and preserve Bayesian coherence under mild regularity conditions
          \citep{grunwald2017inconsistency}.
\end{enumerate}

\begin{algorithm}[H]
\caption{Online Bayesian Update with Tempered Likelihood}
\label{alg:online_update}
\begin{algorithmic}[1]
\Require Current posterior $p = P(H_1 \mid D_{1:n})$, new delta $\Delta$,
         complexity $C$
\Ensure Updated posterior $p' = P(H_1 \mid D_{1:n+1})$
\If{$C \leq 0$}
    \State \Return $p$ \Comment{No evidence from trivial position}
\EndIf
\State $\ell_0 \gets C \cdot \log P(\Delta \mid H_0)$
    \Comment{Tempered human log-lik}
\State $\ell_1 \gets C \cdot \log P(\Delta \mid H_1)$
    \Comment{Tempered cheater log-lik}
\State $\ell_0 \gets \ell_0 + \log(1 - p)$
    \Comment{Add log-prior (human)}
\State $\ell_1 \gets \ell_1 + \log(p)$
    \Comment{Add log-prior (cheater)}
\State $m \gets \max(\ell_0, \ell_1)$
\State $p' \gets e^{\ell_1 - m} \big/ \left(e^{\ell_0 - m} + e^{\ell_1 - m}
    \right)$
\State \Return $p'$
\end{algorithmic}
\end{algorithm}

\subsection{Confidence Estimation}

The confidence in the classification is modulated by two factors:
\begin{equation}
    \text{confidence} = \underbrace{\min\!\left(1, \frac{n}{30}\right)}
    _{\text{data sufficiency}} \times \underbrace{2 \left|P(H_1 \mid D)
    - 0.5\right|}_{\text{posterior certainty}}.
    \label{eq:confidence}
\end{equation}
This ensures that classifications based on few moves are reported with
appropriately low confidence, regardless of how extreme the posterior may be.

\subsection{Formal Likelihood Ratio Test}
\label{sec:lrt}

As a frequentist complement to the Bayesian posterior, we provide a formal
\emph{Likelihood Ratio Test} (LRT) grounded in Neyman--Pearson
theory~\citep{neyman1933most}. The test statistic is:
\begin{equation}
    \Lambda = \frac{P(D \mid H_1)}{P(D \mid H_0)}
    = \exp(L_1 - L_0),
    \label{eq:lrt}
\end{equation}
where $L_0$ and $L_1$ are the log-likelihoods under the human and cheater
hypotheses, respectively. By the Neyman--Pearson lemma, the test that
rejects $H_0$ when $\Lambda > \eta$ is the most powerful test at
significance level $\alpha$ for testing $H_0$ against $H_1$.

The critical value $\eta$ can be calibrated to achieve a desired false
positive rate $\alpha$ (e.g., $\alpha = 0.05$) using the empirical
distribution of $\Lambda$ under $H_0$, obtained from a reference corpus of
confirmed human games. This provides a principled threshold with controlled
error rate, complementing the Bayesian posterior threshold which depends on
the prior.

In log space, the test reduces to comparing:
\begin{equation}
    \log \Lambda = L_1 - L_0 \gtrless \log \eta.
    \label{eq:lrt_log}
\end{equation}

\subsection{Prior Sensitivity Analysis}
\label{sec:sensitivity}

The choice of prior $P(H_1) = 0.01$ is conservative, reflecting a baseline
assumption that only 1\% of players cheat. This choice controls the false
positive rate but makes the posterior slow to converge toward $H_1$ even
under strong evidence.

We recommend conducting a \emph{sensitivity analysis} by evaluating
$P(H_1 \mid D)$ across a range of priors
$P(H_1) \in \{0.001, 0.01, 0.05, 0.10\}$ and reporting the posterior under
each. If the conclusion is robust across priors---i.e., the posterior
exceeds the cheater threshold for all reasonable priors---the evidence is
compelling regardless of the specific prior chosen. Sensitivity analysis
is standard practice in Bayesian applications with subjective
priors~\citep{berger1987testing}.


% =============================================================================
% 6. INFORMATION THEORY
% =============================================================================
\section{Information-Theoretic Metrics}
\label{sec:information_theory}

Complementing the delta-based Bayesian model, we employ information-theoretic
measures to capture distributional properties of player behavior.

\subsection{Shannon Entropy of Move Quality}

The distribution of the player's move quality over a game can be characterized
by its Shannon entropy:
\begin{equation}
    H(P) = -\sum_{i} p_i \log_2 p_i,
    \label{eq:shannon}
\end{equation}
where $p_i$ is the empirical probability of the player's delta falling in bin
$i$ of a histogram over $[0, \max \Delta]$.

Human players exhibit \emph{higher entropy} (more varied move quality), while
engine-assisted players exhibit \emph{lower entropy} (consistently near-optimal
play). This provides a model-free complement to the parametric analysis.

\subsection{KL Divergence from Engine Policy}

For each position, the engine produces a probability distribution over
candidate moves (via softmax of evaluations). The player's actual choice is
represented as a one-hot distribution. The \emph{Kullback--Leibler divergence}
measures how far the player's choice deviates from the engine's recommendation:
\begin{equation}
    D_{\text{KL}}(P_{\text{player}} \| P_{\text{engine}}) = \sum_{i}
    P_{\text{player}}(m_i) \log_2 \frac{P_{\text{player}}(m_i)}
    {P_{\text{engine}}(m_i)},
    \label{eq:kl_divergence}
\end{equation}
with the convention that $P_{\text{engine}}(m_i)$ is floored at $\epsilon =
10^{-10}$ to avoid $\log(0)$.

\begin{remark}
A \emph{low average KL divergence} across a game indicates that the player
consistently selects the moves most favored by the engine---a signature of
engine-assisted play. A high average KL divergence reflects a player who
frequently makes choices the engine would not prioritize, consistent with
human play.
\end{remark}

\begin{remark}[Information Overlap with the Delta Model]
\label{rem:overlap}
Since $P_{\text{player}}$ is a one-hot distribution, the KL divergence
simplifies to $D_{\text{KL}} = -\log_2 Q(m_{\text{chosen}})$, which is
the negative log-likelihood of the chosen move under the engine's policy.
This is informationally related to the winrate delta $\Delta$: both measure
how closely the player's choice aligns with the engine's preference, but in
different spaces (probability space vs.\ winrate space). Practitioners
should be aware that these metrics are \emph{not independent} signals.
The KL divergence is best interpreted as a distributional complement to
$\Delta$, not as a fully orthogonal evidence source.
\end{remark}

\subsection{Cross-Entropy}

The cross-entropy between the player's choice distribution $P$ and the engine's
policy~$Q$ is:
\begin{equation}
    H(P, Q) = -\sum_{i} P(m_i) \log_2 Q(m_i).
    \label{eq:cross_entropy}
\end{equation}
Since $D_{\text{KL}}(P \| Q) = H(P, Q) - H(P)$, the cross-entropy provides
an unnormalized measure of the player's alignment with the engine, combining
the entropy of the player's own distribution with the divergence.


% =============================================================================
% 7. TEMPORAL ANALYSIS
% =============================================================================
\section{Temporal Analysis and Switch-Point Detection}
\label{sec:temporal}

A sophisticated cheater may activate the engine only during critical phases of
a game, playing honestly in the opening before switching to engine-assisted
play when the position becomes complex. The temporal model detects such
behavior.

\subsection{Moving Average Smoothing}

Let $\Delta_1, \Delta_2, \ldots, \Delta_n$ be the sequence of deltas ordered
by move number. The \emph{moving average} with window size $W$ is:
\begin{equation}
    \bar{\Delta}_t = \frac{1}{\min(t, W)} \sum_{j=\max(1,\, t-W+1)}^{t}
    \Delta_j.
    \label{eq:moving_avg}
\end{equation}

\subsection{Switch-Point Detection}

A \emph{switch point} is a move index $t$ where the play quality changes
abruptly. Formally, let:
\begin{align}
    \mu_{\text{before}} &= \text{mean}\!\left(\bar{\Delta}_{t-W}, \ldots,
    \bar{\Delta}_{t-1}\right), \label{eq:mu_before}\\
    \mu_{\text{after}} &= \text{mean}\!\left(\bar{\Delta}_t, \ldots,
    \bar{\Delta}_{t+W-1}\right). \label{eq:mu_after}
\end{align}

A switch point is detected at index~$t$ if:
\begin{equation}
    |\mu_{\text{before}} - \mu_{\text{after}}| \geq \theta_{\text{switch}},
    \label{eq:switch_condition}
\end{equation}
where $\theta_{\text{switch}}$ is the switch threshold (typically 0.10,
representing a 10\% winrate delta change).

\begin{definition}[Switch Direction]
A positive change $\mu_{\text{before}} - \mu_{\text{after}} > 0$ indicates the
player \emph{improved} (deltas decreased---suspicious), while a negative change
indicates \emph{degradation} (deltas increased---natural fatigue).
\end{definition}

Adjacent switch points within a minimum gap of 5~moves are merged, retaining
the one with the largest magnitude, to avoid false triggers from noisy data.

\subsection{Trend and Volatility}

The overall trend in play quality is estimated via simple linear regression:
\begin{equation}
    \hat{\beta} = \frac{\sum_{i=1}^{n}(i - \bar{i})(\Delta_i - \bar{\Delta})}
    {\sum_{i=1}^{n}(i - \bar{i})^2},
    \label{eq:trend_slope}
\end{equation}
where a negative slope indicates improvement over the game.

Volatility is measured by the coefficient of variation:
\begin{equation}
    V = \frac{\sigma_\Delta}{\bar{\Delta}},
    \label{eq:volatility}
\end{equation}
where low volatility may indicate the mechanical consistency of engine-assisted
play.

A game is flagged as \emph{temporally suspicious} if any switch point has
direction ``improved'' and magnitude exceeding $1.5 \times
\theta_{\text{switch}}$.


% =============================================================================
% 8. SYSTEM ARCHITECTURE
% =============================================================================
\section{System Architecture}
\label{sec:architecture}

\subsection{Overview}

The detection system consists of three layers:

\begin{enumerate}
    \item \textbf{Engine Layer}: An external Gomoku engine (e.g., Rapfi)
          communicates via the Gomocup protocol over stdin/stdout.
          For each position, the engine returns the top $k$ candidate moves
          with their winrate evaluations.
    \item \textbf{Analysis Layer}: The mathematical models described in
          Sections~\ref{sec:framework}--\ref{sec:temporal} are applied to
          compute per-move and per-game metrics.
    \item \textbf{Classification Layer}: The Bayesian posterior, information-
          theoretic metrics, and temporal analysis are aggregated into a final
          classification decision.
\end{enumerate}

\subsection{Engine Communication Protocol}

The system sends board positions to the engine using the \texttt{YXBOARD}
command, listing all moves played, followed by \texttt{DONE}. The engine
responds with evaluations. Candidate moves are requested via \texttt{YXNBEST},
which returns the top $k$ moves with their evaluations within a specified time
limit.

After each evaluation cycle, the engine state is reset to ensure independence
between successive queries. This is critical for the independence assumption
underlying the Bayesian model (each $\Delta_i$ is assumed conditionally
independent given the hypothesis).

\subsection{Analysis Pipeline}

For each move in the game, the pipeline performs:
\begin{enumerate}
    \item Send the position (up to the move being analyzed) to the engine.
    \item Receive the top $k$ candidate evaluations.
    \item Compute $\Delta$, complexity $C$, and accuracy $A$.
    \item Perform the tempered Bayesian update.
    \item Compute information-theoretic metrics.
    \item Accumulate temporal statistics.
\end{enumerate}


% =============================================================================
% 9. CLASSIFICATION
% =============================================================================
\section{Classification Decision Framework}
\label{sec:classification}

\subsection{Multi-Signal Aggregation}

The classification integrates three complementary signals:
\begin{enumerate}
    \item \textbf{Bayesian posterior} $P(H_1 \mid D)$ from the tempered
          sequential update (Section~\ref{sec:bayesian}).
    \item \textbf{Boltzmann temperature score} $S_{\tau} = 1/\hat{\tau}$
          from the choice model (Section~\ref{sec:softmax}).
    \item \textbf{Likelihood ratio} $\Lambda$ from the formal LRT
          (Section~\ref{sec:lrt}).
\end{enumerate}

The primary classification is based on the Bayesian posterior:
\begin{equation}
    \text{Classification} = \begin{cases}
        \textbf{Cheater} & \text{if } P(H_1 \mid D) \geq \theta_{\text{cheat}}
            = 0.70, \\
        \textbf{Suspicious} & \text{if } \theta_{\text{susp}} \leq
            P(H_1 \mid D) < \theta_{\text{cheat}}, \\
        \textbf{Human} & \text{if } P(H_1 \mid D) < \theta_{\text{susp}}
            = 0.30.
    \end{cases}
    \label{eq:classification}
\end{equation}

The LRT provides an independent cross-check: if $\log \Lambda > \log \eta$
at significance level $\alpha = 0.05$, the frequentist test also rejects
$H_0$. Agreement between both methods strengthens the conclusion.
The temperature score $S_{\tau}$ serves as corroborating evidence when it
exceeds empirically calibrated thresholds.

\subsection{Summary Statistics}

In addition to the posterior probability, the system reports:
\begin{itemize}
    \item $\hat{\lambda}_{\text{MLE}}$: Estimated rate parameter from the delta
          model.
    \item Mean and standard deviation of $\Delta$.
    \item Near-optimal ratio: fraction of moves with $\Delta < 0.02$.
    \item Boltzmann temperature $\hat{\tau}$ and score $S_{\tau}$.
    \item Log-likelihood ratio $\log \Lambda$.
    \item Average KL divergence from engine policy.
    \item Number and location of temporal switch points.
    \item Confidence score.
\end{itemize}

These supplementary metrics enable human reviewers to understand the basis for
the classification and exercise judgment in borderline cases.


% =============================================================================
% 10. DISCUSSION
% =============================================================================
\section{Discussion}
\label{sec:discussion}

\subsection{Strengths}

\begin{enumerate}
    \item \textbf{Mathematical rigor.} Every component has a clear
          probabilistic interpretation. The tempered likelihood mechanism is
          derived from established Bayesian theory rather than heuristic rules.
    \item \textbf{Complexity awareness.} The framework explicitly accounts for
          position difficulty, reducing false positives from forced sequences
          and false negatives from selective cheaters who avoid using the engine
          in obvious positions.
    \item \textbf{Modularity.} Each model component (delta, complexity,
          Bayesian, information-theoretic, temporal) can be independently
          validated, updated, or replaced.
    \item \textbf{Transparency.} The system provides detailed per-move
          statistics, enabling human review and ensuring accountability.
\end{enumerate}

\subsection{Limitations}

\begin{enumerate}
    \item \textbf{Engine dependency.} The quality of the analysis is bounded by
          the engine's evaluation accuracy. Positions where the engine itself
          errs will produce misleading delta values.
    \item \textbf{Parameter sensitivity.} The default parameters
          ($\lambda_{\text{human}} = 5$, $\lambda_{\text{cheater}} = 50$,
          $\pi = 0.75$, thresholds, etc.) are heuristically chosen. The EM
          algorithm and sensitivity analysis (Sections~\ref{sec:em}
          and~\ref{sec:sensitivity}) partially address this, but full
          calibration requires large labeled datasets.
    \item \textbf{Independence assumption.} The model assumes conditional
          independence of deltas given the hypothesis, which may not hold if
          strategic plans span multiple moves. A Hidden Markov Model
          extension (Section~\ref{sec:conclusion}) can relax this assumption.
    \item \textbf{Selective cheating.} A cheater who uses the engine for only a
          small fraction of critical moves may remain below detection
          thresholds. The temporal model partially addresses this, but
          sufficiently sparse engine use may evade detection.
    \item \textbf{Information overlap.} As noted in
          Remark~\ref{rem:overlap}, the KL divergence and winrate delta are
          not fully independent signals. Over-reliance on correlated metrics
          may lead to overconfident conclusions.
\end{enumerate}

\subsection{Ethical Considerations}

Anti-cheating systems carry inherent risks of false accusations. We advocate
for the following principles:
\begin{enumerate}
    \item Statistical evidence should \emph{inform} human reviewers, not
          replace human judgment.
    \item The threshold for formal action should require high posterior
          probability \emph{and} high confidence (sufficient data).
    \item Players should have access to the evidence and methodology used in
          their case.
    \item The system should be regularly audited for demographic or
          skill-level biases.
\end{enumerate}


% =============================================================================
% 11. CONCLUSION
% =============================================================================
\section{Conclusion and Future Work}
\label{sec:conclusion}

We have presented a comprehensive Bayesian statistical framework for detecting
engine-assisted play in online Gomoku competitions. The framework integrates
parametric move-quality modeling, context-aware complexity weighting, tempered
Bayesian updates, information-theoretic divergence measures, and temporal
switch-point detection into a coherent probabilistic system.

The key technical contributions are:
\begin{enumerate}
    \item A general distributional framework (Exponential, Gamma, Weibull)
          with formal model selection via AIC/BIC, and a two-component
          mixture model with Dirichlet Process extension for human play.
    \item The tempered likelihood mechanism, which provides a mathematically
          principled method for weighting evidence by positional difficulty.
    \item A formal likelihood ratio test grounded in Neyman--Pearson theory
          for controlled false-positive rates, complementing the Bayesian
          posterior.
    \item The integration of multiple complementary detection signals---delta
          distributions, Boltzmann temperature, entropy, KL divergence, and
          temporal patterns---into a unified classification pipeline.
\end{enumerate}

\subsection{Future Work}

Several directions remain for future investigation:
\begin{itemize}
    \item \textbf{Hidden Markov Model for selective cheating:} Instead of
          assuming independence of $\Delta_i$, a two-state HMM with hidden
          states $\{\text{Honest}, \text{Cheating}\}$ would model selective
          engine use as state transitions:
          \begin{equation}
              z_t \mid z_{t-1} \sim \text{Categorical}(A_{z_{t-1}}),
              \quad \Delta_t \mid z_t \sim F_{z_t},
              \label{eq:hmm}
          \end{equation}
          where $A$ is the transition matrix and $F_{z_t}$ is the emission
          distribution for state $z_t$. This subsumes the switch-point
          detection model (Section~\ref{sec:temporal}) as a special case and
          enables probabilistic segmentation of honest vs.\ assisted
          phases~\citep{rabiner1989tutorial}.
    \item \textbf{Empirical calibration:} Fitting all model parameters
          ($\lambda$, $\pi$, $\alpha$, $\beta$, thresholds) to large
          labeled datasets via EM, MCMC, or variational inference.
          Goodness-of-fit tests (Kolmogorov--Smirnov, Anderson--Darling)
          should validate distributional assumptions.
    \item \textbf{Multi-game aggregation:} Extending the sequential Bayesian
          framework to accumulate evidence across multiple games by the same
          player, with a hierarchical prior over player-level parameters.
    \item \textbf{Dirichlet Process deployment:} Replacing the fixed
          two-component mixture with the DPMM (Section~\ref{sec:dpm}) in
          production, allowing the model to automatically discover the
          number of behavioral modes.
    \item \textbf{Adversarial robustness:} Analyzing the system's resilience
          against adaptive cheaters who intentionally introduce errors to
          mimic human play.
    \item \textbf{Generalization:} Adapting the framework to other abstract
          strategy games (e.g., Go, Othello, Connect6) where external engines
          are available.
    \item \textbf{Timing integration:} Incorporating move timing data as an
          additional signal, potentially through a multi-modal Bayesian
          network.
\end{itemize}


% =============================================================================
% REFERENCES
% =============================================================================
\bibliographystyle{plainnat}

\begin{thebibliography}{99}

\bibitem[Regan and Haworth, 2011]{regan2011intrinsic}
Regan, K.~W. and Haworth, G.~M. (2011).
\newblock Intrinsic chess ratings.
\newblock In \emph{Proceedings of the AAAI Conference on Artificial
Intelligence}, volume~25, pages 834--839.

\bibitem[Guid and Bratko, 2006]{guid2006computer}
Guid, M. and Bratko, I. (2006).
\newblock Computer analysis of world chess champions.
\newblock \emph{ICGA Journal}, 29(2):65--73.

\bibitem[Berger and Sellke, 1987]{berger1987testing}
Berger, J.~O. and Sellke, T. (1987).
\newblock Testing a point null hypothesis: The irreconcilability of {P} values
and evidence.
\newblock \emph{Journal of the American Statistical Association},
82(397):112--122.

\bibitem[Gr{\"u}nwald and van Ommen, 2017]{grunwald2017inconsistency}
Gr{\"u}nwald, P. and van Ommen, T. (2017).
\newblock Inconsistency of {B}ayesian inference for misspecified linear models,
and a proposal for repairing it.
\newblock \emph{Bayesian Analysis}, 12(4):1069--1103.

\bibitem[Cover and Thomas, 2006]{cover2006elements}
Cover, T.~M. and Thomas, J.~A. (2006).
\newblock \emph{Elements of Information Theory}.
\newblock Wiley-Interscience, 2nd edition.

\bibitem[Bishop, 2006]{bishop2006pattern}
Bishop, C.~M. (2006).
\newblock \emph{Pattern Recognition and Machine Learning}.
\newblock Springer.

\bibitem[Dempster et~al., 1977]{dempster1977maximum}
Dempster, A.~P., Laird, N.~M., and Rubin, D.~B. (1977).
\newblock Maximum likelihood from incomplete data via the {EM} algorithm.
\newblock \emph{Journal of the Royal Statistical Society: Series B},
39(1):1--38.

\bibitem[Ferguson, 1973]{ferguson1973bayesian}
Ferguson, T.~S. (1973).
\newblock A {B}ayesian analysis of some nonparametric problems.
\newblock \emph{The Annals of Statistics}, 1(2):209--230.

\bibitem[Neal, 2000]{neal2000markov}
Neal, R.~M. (2000).
\newblock Markov chain sampling methods for {D}irichlet process mixture models.
\newblock \emph{Journal of Computational and Graphical Statistics},
9(2):249--265.

\bibitem[Blei and Jordan, 2006]{blei2006variational}
Blei, D.~M. and Jordan, M.~I. (2006).
\newblock Variational inference for {D}irichlet process mixtures.
\newblock \emph{Bayesian Analysis}, 1(1):121--143.

\bibitem[Luce, 1959]{luce1959individual}
Luce, R.~D. (1959).
\newblock \emph{Individual Choice Behavior: A Theoretical Analysis}.
\newblock John Wiley \& Sons.

\bibitem[Neyman and Pearson, 1933]{neyman1933most}
Neyman, J. and Pearson, E.~S. (1933).
\newblock On the problem of the most efficient tests of statistical hypotheses.
\newblock \emph{Philosophical Transactions of the Royal Society of London,
Series A}, 231:289--337.

\bibitem[Rabiner, 1989]{rabiner1989tutorial}
Rabiner, L.~R. (1989).
\newblock A tutorial on hidden {M}arkov models and selected applications in
speech recognition.
\newblock \emph{Proceedings of the IEEE}, 77(2):257--286.

\end{thebibliography}


\end{document}
